% 底本：EGAIV.1.pdf，PDF第180--181页，原分册印刷第181--182页。
\egas{21.9.7}{注记}
（21.9.6，ii）的结论不能推广到 $[k_0:k_0^p]$ 无限的情形。
取 $k_0=\mathbf F_p(X_n)_{n\ge0}$ 为无限多个不定元的有理函数域；
$(X_n)_{n\ge0}$ 是 $k_0/\mathbf F_p$ 的一组 $p$-基，
所以 $[k_0:k_0^p]=\infty$。
\cmt{L12-T08}{底本印刷第181页把这组 $p$-基写成 $(X_n^p)$，这是可从原页确认的印误，而不是识别文字的误差。
每个 $X_n^p$ 已属于 $k_0^p$，其微分为零，不可能构成 $k_0/\mathbf F_p$ 的 $p$-基。
正文改为 $(X_n)$；作者反例的其余构造按原文保留。}
在（21.9.6）中取 $K_0=k_0$，因而只能取 $k=k_0$；
令 $L_0=k_0((T))$，并令 $E=L_0(y)$，其中 $y$ 是多项式
\[
 Y^p-\sum_{n=0}^{\infty}X_nT^{np}\in L_0[Y]
\]
的一个根。
则（21.9.6.2）左、右两端之差 $\delta$ 非零。

事实上，$E/k_0$ 可分；否则 $y^p$ 必属于 $k_0(L_0^p)$，
但由于这里涉及无限多个代数无关的 $X_n$，这是不可能的。
\cmt{L12-G14}{原页这里确实写的是 $y^p$，而不是 $y$；不得把提取文字漏掉上标造成的疑点算成原文错误。
令 $f=y^p=\sum X_n(T^p)^n$。
$k_0(L_0^p)$ 的每个元素只涉及 $k_0$ 中有限多个额外系数，
因而属于某个 $k_0^p(X_i:i\in I)((T^p))$，其中 $I$ 有限；
而 $f$ 的系数含所有 $X_n$，故不属于该域。
由 $p$-基的微分判据，$d_{L_0/k_0}f\ne0$；结合
$E=L_0[Y]/(Y^p-f)$ 的微分表示与 $L_0/k_0$ 可分，
可得 $\Omega^1_{k_0}\otimes_{k_0}E\to\Omega^1_E$ 单射，故 $E/k_0$ 可分（20.6.3）。
这一缺步的完整展开见报告L12-G14。}
所以 $d(E/K_0,k_0)=0$。
（21.9.6.3）化为
\[
 \delta=d(E/M_0,M_0)-1,
 \qquad d(E/M_0,M_0)=\rk_E\Upsilon_{E/M_0}.
\]
只须验证后一个秩为2。
有 $E^p\subseteq M_0$，$E/L_0$ 的 $p$-基由 $y$ 一个元素组成，
而 $L_0/M_0$ 的 $p$-基由 $T$ 一个元素组成。
因为 $E/M_0$ 为代数扩张，由（21.7.1）和（21.4.5），
$\rk_E\Upsilon_{E/M_0}=\rk_E\Omega^1_{E/M_0}=2$。
故 $\delta=2-1=1$，特别地不为零。

\egas{21.9.8}{命题}
设 $k_0$ 为任意特征的域，$K_0/k_0$ 为可分扩张，
$A$ 为一个完备诺特局部 $k_0$-代数，其剩余域 $K$ 是 $K_0$ 的有限扩张。
选定一个与指定剩余域嵌入 $K_0\to K$ 相容的 $k_0$-同态 $K_0\to A$，
并以此赋予 $A$ 及其商环下述 $k$-代数结构；这种提升存在。
\cmt{L12-N08}{原陈述在证明中才选取 $K_0\to A$，但式（21.9.8.1）里的
$\Omega^1_{(A/\mathfrak p)/k}$ 已经需要 $k\to A/\mathfrak p$。
正文先明确这一结构数据，其存在由 $K_0/k_0$ 形式光滑和 $A$ 的完备性给出。
证明适用于任一这样的相容提升；这里不声称这些提升唯一，也不声称它们给出的微分模有额外的典范识别。}
对每个素理想 $\mathfrak p\subset A$，以 $k(\mathfrak p)$ 表示该点的剩余域。
则对每个满足 $k_0\subseteq k\subseteq K_0$、$[K_0:k]<\infty$ 的域 $k$，有
\begin{equation}
\begin{split}
 \rk_{k(\mathfrak p)}
 \bigl(\widehat\Omega^1_{(A/\mathfrak p)/k}
           \otimes_{A/\mathfrak p}k(\mathfrak p)\bigr)
 -\rk_{K_0}\Omega^1_{K_0/k}\\
 \ge \dim(A/\mathfrak p)
       +\rk_{k(\mathfrak p)}\Upsilon_{k(\mathfrak p)/k_0}.
\end{split}\tag{21.9.8.1}
\end{equation}
若此外 $[k_0:k_0^p]<\infty$（$p$ 为 $k_0$ 的特征指数），
则可找到一个域 $k$，满足 $k_0(K_0^p)\subseteq k\subseteq K_0$、$[K_0:k]<\infty$，
使（21.9.8.1）取等号。

\eproof
$A/\mathfrak p$ 与 $A$ 有同一剩余域 $K$，且仍完备，
故可归结为 $A$ 为整环且 $\mathfrak p=0$ 的情形。
置 $E=k(\mathfrak p)=\operatorname{Frac}(A)$。
$K_0/k_0$ 形式光滑（19.6.1），故存在相容的 $k_0$-单同态 $K_0\to A$，
使 $A$ 成为 $K_0$-代数。
又存在 $A$ 的 $K_0$-子代数 $A_0$，它 $K_0$-同构于
$K_0[[T_1,\ldots,T_r]]$，且 $A/A_0$ 有限（19.8.9）。
所以 $E$ 在 $A_0$ 的分式域 $L_0=K_0((T_1,\ldots,T_r))$ 上有限。
由 $K_0/k_0$ 可分，有 $\Upsilon_{K_0/k_0}=0$，并且 $\trdeg_{L_0}E=0$；
由 $A/A_0$ 有限，有 $\dim(A)=\dim(A_0)=r$（16.1.5）。

若 $p=1$，由（21.9.5）以及（17.1.4，iii），
\[
 \rk_E(\widehat\Omega^1_{A/k}\otimes_AE)=r.
\]
$K_0/k$ 为有限可分扩张，故 $\Omega^1_{K_0/k}=0$（21.7.4）；
同时 $E/k_0$ 可分，故 $\Upsilon_{E/k_0}=0$。
于是（21.9.8.1）在这一情形取等号。

现设 $p>1$。
令 $B=k[[T_1,\ldots,T_r]]$，则 $A$ 为有限 $B$-代数。
再令 $B_1=k[[T_1^p,\ldots,T_r^p]]$，
由（21.9.4），$\widehat\Omega^1_{A/k}$ 识别为 $\Omega^1_{A/B_1}$。
若 $M=k((T_1^p,\ldots,T_r^p))$ 为 $B_1$ 的分式域，
则由（20.5.9），$\widehat\Omega^1_{A/k}\otimes_AE$ 识别为 $\Omega^1_{E/M}$。
此外，$\Upsilon_{K_0/k_0}=0$ 和（21.6.1.2）给出
$d(E/K_0,k_0)=\rk_E\Upsilon_{E/k_0}$。
于是（21.9.8.1）就是（21.9.6.2）。
最后，按（21.8.6）的构造选取所需的有限余度子域族，
再应用（21.9.6，ii），即得取等号的 $k$。
\finishproof
